% Part: computability
% Chapter: computability-theory
% Section: ce-closed-cup-cap

\documentclass[../../../include/open-logic-section]{subfiles}

\begin{document}

\olfileid{cmp}{thy}{clo}

\olsection[د c.e. سټونو اتحاد او اشتراک]{په محاسبوي ډول د شمېر وړ
  سټونه د اتحاد او اشتراک لاندې تړلي دي}

لاندې قضيه د محاسبوي ډول د شمېر وړ سټونو د ټولګي ځينې تړلتيايي
خاصيتونه راکوي۔

\begin{thm}
فرض کړئ $A$ او~$B$ په محاسبوي ډول د شمېر وړ دي۔ بيا $A \cap B$ او
$A \cup B$ هم په محاسبوي ډول د شمېر وړ دي۔
\end{thm}

\begin{proof}
\olref[eqc]{thm:ce-equiv} اجازه راکوي چې په محاسبوي ډول د شمېر وړ
سټونو بېلابېلې ځانګړنې وکاروو۔ د روښانتيا دپاره به څو بېل ثبوتونه
وړاندې کړو۔

د لومړي ثبوت دپاره فرض کړئ~$A$ د محاسبه کېدونکې تابع~$f$ په وسيله
شمېرل کېږي، او~$B$ د محاسبه کېدونکې تابع~$g$ په وسيله شمېرل کېږي۔
وټاکئ:
\begin{align*}
h(x) & = \umin{y}{(f(y) = x \lor g(y) = x)} \text{ او}\\
j(x) & = \umin{y}{(f((y)_0) = x \land g((y)_1) = x)}.
\end{align*}
بيا $A \cup B$ د~$h$ د تعريف ساحه، او $A \cap B$ د~$j$ د تعريف
ساحه ده۔

\begin{explain}
په محاسبوي ژبه کښې خبره داسې ده: که داسې کړنلارې راکړل شوې وي چې
$A$ او~$B$ شمېري، نو دا چې يو غړے~$x$ په $A \cup B$ کښې دے که نه،
په دواړو شمېرنو کښې د~$x$ په لټولو نيمه پرېکړه کولے شو؛ او دا چې
$x$ په $A \cap B$ کښې دے که نه، په دواړو شمېرنو کښې په يو وخت د~$x$
په لټولو نيمه پرېکړه کولے شو۔
\end{explain}

د دويم ثبوت دپاره بيا فرض کړئ چې~$A$ د~$f$ او~$B$ د~$g$ په وسيله
شمېرل کېږي۔ وټاکئ:
\[
k(x) = \begin{cases}
f(x/2) & \text{که $x$ جفت وي} \\
g((x-1)/2) & \text{که $x$ طاق وي۔}
\end{cases}
\]
بيا~$k$، $A \cup B$ شمېري؛ نظر دا دے چې~$k$ په نوبت د~$f$ او~$g$
شمېرنې کاروي۔ د~$A \cap B$ شمېرل لږ ستونزمن دي۔ که $A \cap B$ تش
وي، په څرګند ډول په محاسبوي ډول د شمېر وړ دے۔ که نه، $c$ د
$A \cap B$ هر يو غړے وي، او~$l$ داسې تعريف کړئ:
\[
l(x) = \begin{cases}
f((x)_0) & \text{که $f((x)_0) = g((x)_1)$ وي} \\
c & \text{که داسې نۀ وي۔}
\end{cases}
\]
په محاسبوي ژبه کښې~$l$ د~$f$ او~$g$ د شمېرنو د غړو پر جوړو تېرېږي
او هره برابره جوړه چې ومومي، وت ئې راګرځوي؛ که جوړه برابره نۀ وي،
ثابت غړے~$c$ راګرځوي۔
\textbf{د ژباړې د نسخې سمون (OLCMP-027):} سرچينې دې وروستي بديل ته
``د ثابت قيمت په راګرځولو تم کېدل'' ويلي و، خو تابع په هر ننوت يو
قيمت راګرځوي او محاسبه ئې بې پايلې نۀ پاتې کېږي۔ دلته دا د سم ثابت
شاتګ په توګه بيان شوے؛ په همدې برخه کښې دوه څرګندې انګرېزي ليکدوديزې
تېروتنې هم په طبيعي پښتو کښې نۀ دي نقل شوې۔ فورمولونه نۀ دي بدل شوي۔

د وروستي ثبوت دپاره فرض کړئ~$A$ د جزوي تابع~$m(x)$ \emph{د تعريف
ساحه} او~$B$ د جزوي تابع~$n(x)$ د تعريف ساحه ده۔ بيا $A \cap B$ د
جزوي تابع $m(x) + n(x)$ د تعريف ساحه ده۔

\begin{explain}
په محاسبوي ژبه کښې، که~$A$ د هغو قيمتونو سټ وي چې~$m$ پرې درېږي او
$B$ د هغو قيمتونو سټ وي چې~$n$ پرې درېږي، نو $A \cap B$ د هغو
قيمتونو سټ دے چې دواړه کړنلارې پرې درېږي۔
\end{explain}

د درېدنې د قيمتونو د سټ په توګه د~$A \cup B$ څرګندول لږ ستونزمن دي،
ځکه~$m$ او~$n$ بايد په موازي ډول تقليد شي۔ $d$ د~$m$ يو شاخص او~$e$
د~$n$ يو شاخص وي؛ په بله وينا، $m = \cfind{d}$ او
$n = \cfind{e}$۔ بيا $A \cup B$ د دې تابع د تعريف ساحه ده:
\[
p(x) = \umin{y}{(T(d,x,y) \lor T(e,x,y))}.
\]
\begin{explain}
په محاسبوي ژبه کښې، پر ننوت~$x$ تابع~$p$ يا د~$m$ يا د~$n$ درېدلې
محاسبه لټوي، او که يوه ئې ومومي درېږي۔
\end{explain}
\end{proof}

\end{document}
